% Part: conditional-logics
% Chapter: minimal-change-semantics
% Section: true-false

\documentclass[../../../include/open-logic-section]{subfiles}

\begin{document}

\olfileid{con}{min}{tf}

\olsection{પ્રતિતથ્યાત્મક વિધાનોની સત્યતા અને અસત્યતા}

સૌથી નજીકનાં બધાં $!A$-જગતો $!B$-જગતો હોય, તો
$!A \cif !B$ પ્રતિતથ્યાત્મક વિધાન રિક્તતા વગર સત્ય હોય છે;
તે~\olref{fig:true}માં દર્શાવ્યું છે.
\begin{figure}
\begin{center}
\begin{tikzpicture}[scale=.7]\small
  \spheresystem{5}
  \draw (0,0) node {$w$};
  \propositionintersect{0}{4}{40}{3.5}
  {\tikzset{proposition/.append={smooth,tension=1.4}}
  \proposition[shift={(1.6,-1)}]{90}{1}{30}{4}}
  \path (1.5,3) node[below] {$\formula{B}$};
  \path (3,0) node {$\formula{A}$};
\end{tikzpicture}
\caption{રિક્તતા વગર સત્ય પ્રતિતથ્યાત્મક વિધાન}
\ollabel{fig:true}
\end{center}
\end{figure}
$w$ની આસપાસના ગોળાઓના તંત્રમાં $!A$ને સ્વીકારતો એક પણ ગોળો ન
હોય તોપણ પ્રતિતથ્યાત્મક વિધાન~$w$ પર સત્ય છે. એવા કિસ્સામાં તે
\emph{રિક્તતાને કારણે} સત્ય છે (જુઓ~\olref{fig:vacuous}).
\begin{figure}
\begin{center}
\begin{tikzpicture}[scale=.7]\small
  \spheresystem{5}
  \draw (0,0) node {$w$};
  \proposition{0}{6.5}{40}{3.5}
  {\tikzset{proposition/.append={smooth,tension=1.4}}
  \proposition[shift={(1.2,-1)}]{90}{1}{30}{4}}
  \path (1.5,3) node[below] {$\formula{B}$};
  \spherepos{0}{7}{node {$\formula{A}$}}
\end{tikzpicture}
\caption{રિક્તતાને કારણે સત્ય પ્રતિતથ્યાત્મક વિધાન}
\ollabel{fig:vacuous}
\end{center}
\end{figure}

તે બે રીતે અસત્ય હોઈ શકે છે. પહેલી રીતે, સૌથી નજીકનાં $!A$-જગતોમાં
કેટલાક $!B$-જગતો હોય પરંતુ બધાં ન હોય. આ કિસ્સામાં
$!A \cif \lnot !B$ પણ અસત્ય છે (જુઓ~\olref{fig:false}).
\begin{figure}
\begin{center}
\begin{tikzpicture}[scale=.7]\small
  \spheresystem{5}
  \draw (0,0) node {$w$};
  \propositionintersect{0}{4}{40}{3.5}
  {\tikzset{proposition/.append={smooth,tension=1.4}}
  \proposition[shift={(1.6,0)}]{90}{1}{30}{3}}
  \path (1.5,3) node[below] {$\formula{B}$};
  \path (3,0) node {$\formula{A}$};
\end{tikzpicture}
\caption{અસત્ય પ્રતિતથ્યાત્મક વિધાન; વિરુદ્ધ પરિણામવાળું પણ અસત્ય}
\ollabel{fig:false}
\end{center}
\end{figure}
સૌથી નજીકનાં $!A$-જગતો અને $!B$-જગતો એકબીજાને જરા પણ ન મળે,
તો $!A \cif !B$ અસત્ય છે. પરંતુ આ કિસ્સામાં સૌથી નજીકનાં બધાં
$!A$-જગતો $\lnot !B$-જગતો છે; તેથી $!A \cif \lnot !B$ સત્ય છે
(જુઓ~\olref{fig:false-opposite}).
\begin{figure}
\begin{center}
\begin{tikzpicture}[scale=.7]\small
  \spheresystem{5}
  \draw (0,0) node {$w$};
  \propositionintersect{0}{4}{40}{3.5}
  {\tikzset{proposition/.append={smooth,tension=1.4}}
  \proposition[shift={(-1,0)}]{90}{1}{30}{3}}
  \path (-1,3) node[below] {$\formula{B}$};
  \path (3,0) node {$\formula{A}$};
\end{tikzpicture}
\caption{અસત્ય પ્રતિતથ્યાત્મક વિધાન; વિરુદ્ધ પરિણામવાળું સત્ય}
\ollabel{fig:false-opposite}
\end{center}
\end{figure}

કડક શરતી વિધાનથી વિપરીત, પ્રતિતથ્યાત્મક વિધાનો સંજોગાધીન હોઈ શકે
છે. \olref{fig:contingent}માંનો ગોળા-નિદર્શ જુઓ. $u$થી સૌથી નજીકનાં
$!A$-જગતો બધાં $!B$-જગતો છે, તેથી $\mSat{M}{!A \cif
  !B}[u]$. પરંતુ~$v$થી સૌથી નજીકનાં કેટલાક $!A$-જગતો $!B$-જગતો
નથી, તેથી $\mSat/{M}{!A \cif !B}[v]$.
\begin{figure}
\begin{center}
\begin{tikzpicture}[scale=.6]\small
  \spheresystem{7}
  \spheresystem[shift={(0:3.1)},dashed]{4}
  \propositionintersect{45}{4}{40}{4.5}
  \begin{scope}
    \clip \propositionplot{45}{4}{40}{4.5} ;
    \spherelayer[shift={(0:3.1)}]{3}
  \end{scope}
  \proposition[shift={(1.4,-.5)}]{90}{1}{35}{5}
  \path (0,0) node {$u$};
  \path (0:3.1) node {$v$};
  \spherepos{35}{9}{node {$\formula{A}$}}
  \spherepos{80}{9}{node {$\formula{B}$}}
\end{tikzpicture}
\end{center}
\caption{સંજોગાધીન પ્રતિતથ્યાત્મક વિધાન}
\ollabel{fig:contingent}
\end{figure}
\end{document}
